\section{Related Work}
\label{sec:related_work}

\subsection{Vision-Language-Action Models}
Generalist robot manipulation has advanced substantially with the convergence of vision-language foundation models and robotic control \cite{open_x_embodiment_rt_x_2023,kim2024openvla,black2024pi_0,octo2023}. Early models framed manipulation as autoregressive token generation, quantizing continuous end-effector actions into discrete language vocabulary bins (e.g., RT-1, RT-2 \cite{open_x_embodiment_rt_x_2023}, OpenVLA \cite{kim2024openvla}). While discrete tokenization enables direct transfer from pre-trained Large Language Models, it suffers from discretization errors and high sequence generation latency.

To achieve fine-grained, continuous multi-modal control, modern architectures employ generative diffusion and conditional flow matching \cite{chi2023diffusion,peebles2023scalable,black2024pi_0,shukor2025smolvla}. Policies such as Octo \cite{octo2023} and $\pi_0$ \cite{black2024pi_0} denoise continuous action trajectories conditioned on vision and language features. However, existing generalist models typically scale from 3B to over 10B parameters, requiring enterprise-grade GPU clusters ($>$16--40\,GB VRAM) and hindering practical deployment on resource-constrained physical robots. Recent efforts toward compact policies (e.g., SmolVLA \cite{shukor2025smolvla}, RoboMamba \cite{liu2024robomamba}) explore parameter reduction; however, they remain reliant on 2D image coordinates, leaving spatial robustness unresolved.

\subsection{3D Geometry in Robot Visuomotor Learning}
Standard VLA architectures treat visual observations as planar 2D grids using learnable 2D positional embeddings or standard sine-cosine encodings. While effective for semantic scene parsing, 2D representations discard crucial physical camera geometry, rendering policies brittle to camera extrinsics variations \cite{liberomax2024}. Explicit 3D representations, such as 3D point clouds or voxel grids, directly represent physical coordinates but depend on active depth sensors (which fail on reflective, transparent, or dark surfaces) and introduce heavy computational overhead during voxelization and 3D convolutions.

Rotary Position Embeddings (RoPE) \cite{su2024roformer} have demonstrated remarkable length generalization in transformer models. In robotics, ray-based positional encodings have recently emerged to ground RGB observations in physical 3D space by mapping each image patch to a 3D ray direction vector originating from the optical center. Unlike full point cloud pipelines, continuous ray-based embeddings operate directly on standard RGB inputs while incorporating camera intrinsic and extrinsic transformations $[\mathbf{R} \mid \mathbf{t}]$. Our work extends this formulation to a continuous 3D Ray-RoPE mechanism within a flow-matching action expert, achieving metric coordinate invariance under viewpoint shifts without requiring active depth sensors.

\subsection{Action Chunking and Kinematic Smoothness}
Action chunking—predicting a continuous temporal trajectory of future actions $[\mathbf{a}_t, \mathbf{a}_{t+1}, \dots, \mathbf{a}_{t+H}]$ at each control step—has become standard in visuomotor control (e.g., ACT \cite{zhao2023act}, Diffusion Policy \cite{chi2023diffusion}). Chunking prevents the policy from compounding single-step autoregressive drift. However, action chunking introduces the \textit{chunk-boundary discontinuity problem}: because each temporal horizon is generated independently through stochastic diffusion denoising, adjacent chunks exhibit high velocity and acceleration discrepancies at boundary transitions.

Standard practitioners mitigate this discontinuity using heuristic temporal moving averages or exponential weighting across overlapping predictions \cite{zhao2023act}. While temporal averaging smooths boundaries, it introduces phase lag, increases inference latency by executing redundant denoising steps, and can cause trajectory flattening near rapid directional reversals. In contrast, VGA explicitly conditions the flow-matching denoiser on the preceding executed trajectory ($\mathbf{a}_{\text{prev}}$) and incorporates an analytical boundary velocity penalty ($\mathcal{L}_{\text{smooth}}$), guaranteeing $C^1$ kinematic continuity directly during trajectory generation.

% [PLACEHOLDER / LITERATURE NOTE: Additional related work on state-space models (e.g., Mamba-based robotic policies \cite{gu2023mamba,liu2024robomamba}) or specific tactile/force-feedback integration may be cited here if evaluated in follow-up revisions.]
